\section{Practical Value of Directional Adaptation}
\label{sec:discussion}
\paragraph{Better ranking without a larger inference model.}
\method\ combines the strongest development macro quality in the evaluated
panel with a compact deployment artifact. For 384-to-768-dimensional
conversion, its 295,680 float32 coefficients occupy 1.13\,MiB: 114.7 times
fewer parameter bytes than the evaluated 33,927,168-parameter Converter.
The two corrections are absorbed into the affine coefficients, so they add
no stages to document mapping at inference. Affine ridge has the same
coefficient count; the comparison therefore separates the retrieval
benefit of the correction procedure from the storage benefit of affine
mapping relative to the evaluated neural alternative.

\paragraph{A clear operating scenario.}
The intended user must support destination-model queries while retaining
a single destination-space gallery during staged document replacement.
The fixed-anchor study demonstrates strong mapping quality with limited
paired supervision; the all-paid MiniLM result shows a useful operating
point when all available paid pairs are used. Prospective E5 gains over
Converter and ridge persist through 25\%, 50\% and 75\% coverage. These
gains complement the TREC-COVID/MiniLM operating point: $0.5195$ nDCG at
75\% coverage versus $0.5203$ for native Full-Y. These measured conditions
make \method\ a concrete candidate for evaluating a required
encoder transition. If source queries can remain in service, native
Full-X remains a viable alternative: every partial development five-seed
ARC2 condition mean is below it. Procrustes and LA2M also remain strong choices on
the prospective collection, so encoder-specific mixed-gallery validation
should determine which adapter to deploy.

\paragraph{What the efficiency evidence measures.}
The compact final map and absence of neural mapper training are structural
advantages. Our cached-vector experiments measure retrieval under matched
destination exposure. Physical encoding time, fitting workspace, index
resources and serving SLOs are separate deployment quantities; the
parameter-storage ratio is not an end-to-end cost ratio.
